% from ega4-2-fr.tex
\oldpage[IV-2]{84}

un corps de caractéristique $p>0$, $k$ le corps des fractions
rationnelles $k_0(s,t)$ à $2$ indéterminées, $L$ le corps
$k(s^{1/p},t^{1/p})$ et $X=\operatorname{Spec}(L)$. Soient d'autre part $u$
une troisième indéterminée et soit $K$ le corps
$k(u,s^{1/p}+ut^{1/p})$; on vérifie aisément que $k$ est algébriquement fermé
dans $K$ et que $L\otimes_kK$ a un nilradical non nul, donc $X_{(K)}$ n'est
pas réduit; mais comme $(X_{(K)})_{\mathrm{red}}$ n'est pas défini sur $k$,
et que $K$ ne contient aucune sous-extension de degré fini sur $k$ autre que
$k$, le plus petit corps de définition de $(X_{(K)})_{\mathrm{red}}$ ne peut
être fini sur $k$.

\item[(ii)] Étant donné un élément de l'un des ensembles d'arrivée des
applications (4.8.2), et une extension $K_1$ de $K$, il est équivalent de dire
que l'objet considéré admet un plus petit corps de définition ou que son image
par $\varepsilon_{K_1,K}$ admet un plus petit sous-corps de définition (les
deux sous-corps étant nécessairement les mêmes), comme on l'a vu dans
(4.8.4).
\end{enumerate}
\end{remarks}
\subsection{Corps de définition d'une partie d'un préschéma}

\begin{proposition}[4.9.1]
\label{IV.2.4.9.1-fr}
Soient $k$ un corps, $X$ un $k$-préschéma, $K$ une extension de $k$,
$\overline T$ une partie fermée de $X_{(K)}$, $K'$ une sous-extension de $K$.
Les conditions suivantes sont équivalentes:
\begin{enumerate}
\item[a)] $\overline T$ est définie sur $K'$.
\item[b)] La partie ouverte $\overline U=X_{(K)}-\overline T$ de $X_{(K)}$
est définie sur $K'$.
\item[b')] Le sous-préschéma de $X_{(K)}$ induit sur l'ouvert $\overline U$
est défini sur $K'$.
\item[c)] Il existe un sous-préschéma fermé $\overline Z$ de $X_{(K)}$ ayant
$\overline T$ pour espace sous-jacent, et défini sur $K'$.
\end{enumerate}
\end{proposition}

\begin{proof}
Si $g:X_{(K)}\to X_{(K')}$ est la projection canonique, dire que
$\overline T$ est définie sur $K'$ signifie qu'il existe une partie fermée
$T'$ de $X_{(K')}$ telle que $\overline T=g^{-1}(T')$; comme
$g^{-1}(X_{(K')}-T')=X_{(K)}-g^{-1}(T')$ pour toute partie $T'$ de
$X_{(K')}$, a), b) et b') sont équivalentes. Si $\overline Z$, ayant
$\overline T$ pour espace sous-jacent, est défini sur $K'$, on a
$\overline Z=Z'\otimes_{K'}K$, où $Z'$ est un sous-préschéma fermé de
$X_{(K')}$, et si $T'$ est le sous-espace sous-jacent à $Z'$, on sait que
$\overline T=g^{-1}(T')$ (I, 4.4.1), donc $\overline T$ est définie sur
$K'$. Inversement, pour voir que a) entraîne c), il suffit de considérer un
sous-préschéma fermé $Z'$ de $X_{(K')}$ ayant pour espace sous-jacent $T'$
(I, 5.2.1) et de prendre $\overline Z=Z'\otimes_{K'}K$.
\end{proof}

\begin{corollary}[4.9.2]
\label{IV.2.4.9.2-fr}
Avec les notations de (4.9.1), supposons que $K'$ soit un corps de définition
de $\overline T$; toute extension $K''\subset K'$ de $k$, telle que $K'$ soit
extension radicielle de $K''$, est aussi un corps de définition de
$\overline T$.
\end{corollary}

\begin{proof}
Il suffit d'observer que la projection canonique
$g:X_{(K')}\to X_{(K'')}$ est un morphisme entier, surjectif et radiciel,
donc un homéomorphisme (2.4.5).
\end{proof}

\begin{remark}[4.9.3]
\label{IV.2.4.9.3-fr}
Il résulte de (4.9.2) qu'une partie fermée $\overline T$ de $X_{(K)}$ n'admet
pas nécessairement de \emph{plus petit} corps de définition. Par exemple, si
$K=k(t)$, où $t$ est une indéterminée et $k$ un corps parfait de
caractéristique $p>0$, une partie fermée $\overline F$ de $X_{(K)}$ ne peut
avoir de plus petit corps de définition que si elle est déjà définie sur $k$;
en effet, pour tout sous-corps $K'$ de $K$ contenant $k$ et $\ne k$, on a
$K'^p\ne K'$, $K'^p$ contient $k$ et $K'$

